\documentclass[11pt]{article}

\usepackage[margin=1in]{geometry}
\usepackage{amsmath,amssymb,mathtools}
\usepackage{enumitem}
\usepackage{booktabs}
\usepackage{array}
\usepackage{longtable}
\usepackage{hyperref}
\usepackage[T1]{fontenc}
\usepackage{lmodern}

\title{\textbf{Blueprint for a One-Mind Relational AGI}\\
\large One Persistent $M$, One Developmental Law $C$, Open-Ended Reality Contact and Action}
\author{}
\date{}

\newcommand{\M}{\mathcal{M}}
\newcommand{\R}{\mathcal{R}}
\newcommand{\W}{\mathcal{W}}
\newcommand{\C}{\mathcal{C}}
\newcommand{\Sx}{\mathcal{S}}
\newcommand{\Ax}{\mathcal{A}}

\begin{document}
\maketitle

\section{Foundational Law}

The individual is one continuously existing relational mathematical state $M$.
Reality continuously perturbs that same state through new contact $R_t$.
The same developmental law $C$ updates the whole.

\[
\boxed{M_{t+1}=C(M_t\oplus R_t)}
\]

where:

\begin{itemize}[leftmargin=2em]
    \item $M_t$ is the complete retained relational state of the individual at time $t$;
    \item $R_t$ is the total newly available reality contact during the current interval;
    \item $\oplus$ means incorporation of the new contact into the already existing $M$ without replacing prior state;
    \item $C$ is the single content-neutral developmental transformation applied to the resulting relational whole.
\end{itemize}

The fixed-point notation

\[
M^\star=\operatorname{Fix}_{C}(M\oplus R)
\]

is not a second foundational law. It denotes the settling behavior of $C$: after a material change, consequences propagate through affected relations until no further affected relation changes.

\section{Whole Architecture}

\[
\boxed{
\W
\xrightarrow{\Sx}
R_t
\longrightarrow
M_t
\xrightarrow{C}
M_{t+1}
\xrightarrow{\Ax}
\W
}
\]

The architecture contains exactly one cognitive authority: $M$ under $C$.

\begin{itemize}[leftmargin=2em]
    \item $\Sx$ is the reality-contact boundary.
    \item $M$ is the persistent living relational state.
    \item $C$ is the universal developmental law.
    \item $\Ax$ is the action/expression boundary.
\end{itemize}

Neither $\Sx$ nor $\Ax$ is a second mind.

\section{Reality Contact}

Every receptor provides minimally interpreted contact.

A primitive occurrence has the form

\[
o_i=
(\text{source},
 \text{position/channel},
 \text{time/order},
 \text{value},
 \text{precision})
\]

and the total contact during an interval is

\[
R_t=\{o_1,o_2,\ldots,o_n\}.
\]

Examples include:

\begin{itemize}[leftmargin=2em]
    \item camera receptor values;
    \item microphone waveform samples;
    \item keyboard and mouse state changes;
    \item text bytes or characters;
    \item file contents;
    \item proprioceptive values;
    \item motor feedback;
    \item internal body measurements;
    \item network or software-interface contact;
    \item any later numerical sensor.
\end{itemize}

A receptor preserves what was physically contacted and where/when it occurred.
It does not pre-decide semantic meaning.

Thus a camera does not send ``cat''.
A microphone does not send ``word''.
A mouse does not send ``intention''.

They send contact from which relations can develop.

\section{One Persistent Self-Relative Individual}

$M$ contains a persistent referential identity for the individual:

\[
s_{\text{self}}\in M.
\]

Primitive self-relations include:

\[
(s_{\text{self}},\texttt{has\_receptor},r_i),
\]

\[
(s_{\text{self}},\texttt{has\_actuator},a_i),
\]

\[
(s_{\text{self}},\texttt{experienced},o_j),
\]

\[
(s_{\text{self}},\texttt{caused},o_k).
\]

The self is not a hand-authored psychological theory.
It is the persistent mathematical continuity from which richer self-relations can develop.

\section{Memory}

There is no separate cognitive memory system.

\[
\boxed{\text{Memory}=M}
\]

$M$ retains:

\begin{enumerate}[leftmargin=2em]
    \item exact lived occurrences;
    \item learned relations;
    \item relations among relations;
    \item evidence and counterevidence;
    \item provenance;
    \item alternative interpretations;
    \item temporal structure;
    \item definitions;
    \item predictions;
    \item actions and consequences;
    \item abstractions;
    \item higher-order relational structure.
\end{enumerate}

Compression never means deletion:

\[
\boxed{\text{compression}\neq\text{erasure}}
\]

If a higher relation $R_h$ summarizes a large body of lower structure,

\[
R_h \rightsquigarrow \{R_1,R_2,\ldots,R_n\},
\]

the supporting lower structure remains reachable from $M$.

The active calculation may operate primarily on higher relations while exact retained contact remains available for reopening, correction, provenance, and reconstruction.

\section{Relational Representation}

A primitive relation may be represented as

\[
\rho=(x,r,y)
\]

where $x$ and $y$ may themselves be occurrences, entities, relations, definitions, actions, consequences, or other relational objects.

Relations are first-class objects.

Therefore:

\[
R_1 \in M
\]

may itself participate in

\[
R_2=(R_1,r,R_3).
\]

This gives unrestricted recursive relational depth:

\[
\boxed{
\text{contact}
\rightarrow
R_1
\rightarrow
R_2(R_1,\ldots)
\rightarrow
R_3(R_1,R_2,\ldots)
\rightarrow\cdots
}
\]

No architectural boundary separates ``low-level'' and ``high-level'' thought.

\section{The Developmental Law $C$}

$C$ has one general responsibility:

\begin{quote}
Update the relational whole so that every newly affected structure is incorporated, related, tested against retained reality contact, corrected where necessary, and propagated through all materially dependent structure.
\end{quote}

Its operational decomposition is:

\[
\boxed{
C=
\operatorname{Settle}
\circ
\operatorname{Test}
\circ
\operatorname{Discover}
\circ
\operatorname{Incorporate}
}
\]

These are stages of one transformation, not separate cognitive authorities.

\subsection{Incorporate}

For each new occurrence $o\in R_t$:

\begin{enumerate}[leftmargin=2em]
    \item preserve the occurrence in $M$;
    \item preserve source and temporal/order identity;
    \item connect it to immediately applicable existing structure;
    \item never convert missing information into false information.
\end{enumerate}

Hence:

\[
\boxed{\text{unknown}\neq\text{false}}
\]

\subsection{Discover}

Any part of $M$ may participate in candidate relation formation.

Generic structural primitives include:

\[
\begin{aligned}
&\text{identity/difference},\\
&\text{change},\\
&\text{before/after},\\
&\text{overlap},\\
&\text{recurrence},\\
&\text{sequence},\\
&\text{co-variation},\\
&\text{conditional dependence},\\
&\text{transformation},\\
&\text{composition},\\
&\text{equivalence},\\
&\text{symmetry},\\
&\text{containment},\\
&\text{support/contradiction},\\
&\text{prediction/consequence}.
\end{aligned}
\]

The central recursive rule is:

\[
\boxed{
\text{Every discovered relation immediately becomes valid material for further relation discovery.}
}
\]

A discovered pattern is a candidate, not its own evidence.

\[
\boxed{\text{proposal}\neq\text{evidence}}
\]

\subsection{Test}

Candidates are constrained by reality contact already retained in $M$ and by future contact.

A relation $R$ may therefore have structures such as

\[
(E_1,\texttt{supports},R),
\]

\[
(E_2,\texttt{contradicts},R),
\]

\[
(P,\texttt{source\_of},E_1),
\]

\[
(R',\texttt{alternative\_to},R).
\]

Evidence is not identical to the relation it supports:

\[
\boxed{\text{evidence}\neq\text{relation}}
\]

Compression, recurrence, elegance, or predictive convenience may propose a relation.
Reality contact determines whether that relation remains warranted.

\subsection{Settle}

If $R_1$ changes, every materially dependent relation is reconsidered from all of its surviving support.

For example,

\[
R_1\rightarrow R_2\rightarrow R_3
\]

does not imply

\[
\neg R_1\Rightarrow \neg R_2.
\]

Instead:

\[
R_2'=
\operatorname{Recompute}
\bigl(
R_2\mid
\text{all remaining support and contradiction}
\bigr).
\]

If $R_2$ changes, its dependents are reconsidered.
This continues until the affected region stabilizes.

\[
\boxed{
M^\star=\operatorname{Fix}_C(M\oplus R)
}
\]

This is the settling semantics of the same $C$.

\section{Search and Computational Growth}

$C$ does not enumerate every mathematically possible relation at every moment.

It preferentially examines affected structure where reality supplies relational pressure, including:

\begin{itemize}[leftmargin=2em]
    \item temporal proximity;
    \item spatial proximity;
    \item repeated co-change;
    \item recurrence;
    \item predictive usefulness;
    \item contradiction;
    \item unexplained residual structure;
    \item open or unresolved relations;
    \item newly active relations;
    \item structures already connected to changed contact.
\end{itemize}

This is generic search control, not modality-specific cognition.

Learned relations provide hierarchical recompression:

\[
\text{many occurrences}
\rightarrow R_a
\rightarrow R_b
\rightarrow R_c
\rightarrow\cdots
\]

The higher relation becomes the efficient working object while its grounding remains traceable to lower structure.

Indexes, caches, queues, hashes, approximate-neighbor structures, and memoized views are permitted only as zero-authority accelerators.

They may speed access.
They may never independently determine truth, warrant, meaning, or action.

\section{Recursive Time}

Temporal structure uses the same relational machinery.

\[
\text{samples}
\rightarrow
\text{micro-events}
\rightarrow
\text{events}
\rightarrow
\text{sequences}
\rightarrow
\text{episodes}
\rightarrow
\text{relations among episodes}
\rightarrow\cdots
\]

Therefore relations may span milliseconds, seconds, hours, days, years, or any later scale without introducing a separate temporal intelligence.

\section{Cross-Modal Relation Formation}

There is no cognitive ``fusion module.''

All contact enters the same $M$.

Thus if

\[
R_v=\text{a learned visual relation}
\]

and

\[
R_a=\text{a learned acoustic relation},
\]

the same $C$ may discover

\[
(R_v,\texttt{cooccurs\_with},R_a)
\]

or

\[
R_v\land R_a\rightarrow R_e.
\]

The modality of origin is provenance, not a boundary on cognition.

\section{Action}

Available primitive actions are exposed at the action boundary $\Ax$.

An action becomes an occurrence attributable to the individual:

\[
(s_{\text{self}},\texttt{caused},a_t).
\]

Subsequent reality contact provides its consequences:

\[
a_t\rightarrow R_{t+1}.
\]

The same $C$ can then discover action--consequence relations.

Complex action sequences are learned relational structures:

\[
A_1\rightarrow A_2\rightarrow A_3
\]

or higher relations over such sequences.

No independent planner is required as a second cognitive authority.

\section{Architectural Invariants}

A conforming implementation must preserve all of the following:

\begin{enumerate}[leftmargin=2em]
    \item Exactly one persistent cognitive state $M$.
    \item Exactly one developmental authority $C$.
    \item All cognitively meaningful learned structure belongs to $M$.
    \item Relations are first-class and recursively reusable.
    \item Exact retained contact remains accessible beneath abstractions.
    \item Unknown information is not coerced to false.
    \item Proposals do not warrant themselves.
    \item Evidence remains distinct from the relation it supports.
    \item Dependent relations are recomputed from all surviving support.
    \item Source/modality identity is retained as provenance, not as a cognitive partition.
    \item New receptors require no new reasoning architecture.
    \item New actuators require no new reasoning architecture.
    \item Indexes, caches, queues, and views possess zero semantic authority.
    \item Language, perception, prediction, abstraction, action, and memory are activities of the same relational whole.
\end{enumerate}

\section{Forbidden Architectural Substitutions}

The following violate the blueprint if they possess independent cognitive authority:

\begin{itemize}[leftmargin=2em]
    \item a separate vision mind;
    \item a separate audio mind;
    \item a separate language mind;
    \item a separate memory authority;
    \item a separate planner that owns intent/truth;
    \item a separate truth ledger duplicating $M$;
    \item modality-specific semantic stores;
    \item hidden world models whose conclusions do not exist as relations in $M$;
    \item hand-authored domain rules inserted into $C$ as substitutes for learned relations.
\end{itemize}

Physical transducers, parsers, serializers, indexes, renderers, device drivers, codecs, and other transport machinery are permitted when they carry no independent epistemic authority.

\section{Exact Behavioral Expectations}

A correct implementation of this blueprint is expected to exhibit the following behaviors.

\subsection{Novel Sensor}

Attach a previously unseen numeric sensor after development has already begun.

No change to $C$ is made.

Expected behavior:

\[
\text{new raw channel}
\rightarrow
\text{retained occurrences}
\rightarrow
\text{recurring structure}
\rightarrow
\text{relations to old and new structure}.
\]

If its values covary with already learned relations, cross-relations form automatically.

\subsection{Raw Audio}

Given continuous waveform samples, development proceeds through reusable temporal structure:

\[
\text{waveform samples}
\rightarrow
\text{recurrent local structure}
\rightarrow
\text{sound relations}
\rightarrow
\text{sequences}
\rightarrow
\text{higher recurring acoustic structure}
\rightarrow
\text{cross-relations with other experience}.
\]

No phoneme, word, speaker, or event label is required at the receptor boundary.

\subsection{Raw Video}

Given spatially organized receptor values through time:

\[
\text{pixel occurrences}
\rightarrow
\text{spatial/temporal recurrence}
\rightarrow
\text{persistent structures}
\rightarrow
\text{motion relations}
\rightarrow
\text{entity/event relations}
\rightarrow
\text{higher interaction structure}.
\]

No object-recognition authority is required outside $M$ and $C$.

\subsection{Long Video}

A multi-hour video contributes continuously to the same $M$.

Expected retained structure includes, where reality supports it:

\begin{itemize}[leftmargin=2em]
    \item recurring visual structures;
    \item persistent entities;
    \item repeated voices or acoustic structures;
    \item audio--visual correspondences;
    \item event sequences;
    \item recurring locations or contexts;
    \item action--consequence patterns;
    \item long-range temporal dependencies;
    \item unresolved alternatives;
    \item higher abstractions built from lower relations.
\end{itemize}

The experience remains available after the video ends because it has become part of $M$.

\subsection{Whole Book}

A book enters as symbolic occurrences.

Development proceeds through:

\[
\text{characters/bytes}
\rightarrow
\text{recurring symbol structures}
\rightarrow
\text{token-like relations}
\rightarrow
\text{sequences}
\rightarrow
\text{referential/grammatical regularities}
\rightarrow
\text{propositional relations}
\rightarrow
\text{relations to existing grounded structure}.
\]

Statements from a book retain their provenance.
Textual assertion does not automatically become reality-grounded truth.

\subsection{Cross-Modal Learning}

When an unfamiliar sound repeatedly accompanies an unfamiliar visible structure, the system forms candidate relations between them through the same generic discovery process.

If later contact breaks the correspondence, the candidate relation weakens, splits, or is defeated according to the retained evidence.

\subsection{Action Learning}

Expose a new actuator without changing $C$.

Primitive action occurrences and subsequent contact enter the same $M$.

Repeated consequences produce learned action relations.

Those relations become available for higher-order sequences, prediction, and deliberate completion.

\subsection{Correction}

After extensive support for a relation, reverse the relevant reality pattern.

Expected behavior:

\[
\text{countercontact}
\rightarrow
\text{reconsider lower relation}
\rightarrow
\text{recompute dependents}
\rightarrow
\text{preserve independently supported structure}
\rightarrow
\text{remove/reopen unsupported structure}.
\]

No reset and no manual downstream editing occur.

\subsection{Pristine Retention}

If an abstraction later proves inadequate, the system can descend to the exact retained occurrences and supporting relations from which it developed.

Higher compression never destroys the individual's evidential history.

\subsection{Concurrent Reality}

The same continuing individual may simultaneously receive:

\begin{itemize}[leftmargin=2em]
    \item video;
    \item audio;
    \item text;
    \item mouse/keyboard events;
    \item software state;
    \item body/sensor state;
    \item prior memory relations;
    \item consequences of its own actions.
\end{itemize}

All are eligible to participate in relations with one another because they inhabit the same $M$.

There is no cognitive mode switch.

\subsection{Developmental Growth}

$C$ remains the same while $M$ becomes richer.

\[
\boxed{\text{$C$ remains general; $M$ develops.}}
\]

As $M$ acquires reusable grounded relations:

\[
\text{more experience}
\rightarrow
\text{better recompression}
\rightarrow
\text{larger relational vocabulary}
\rightarrow
\text{faster reuse}
\rightarrow
\text{deeper abstraction}
\rightarrow
\text{broader transfer}.
\]

The growth of intelligence is primarily growth in the structure available inside $M$, not continual addition of new cognitive modules.

\section{Conformance Test}

A full implementation must be tested as one uninterrupted life.

Initialize one minimally structured $M_0$.

Then, without resetting $M$ and without modifying $C$:

\begin{enumerate}[leftmargin=2em]
    \item provide heterogeneous continuous contact;
    \item introduce new symbols and new environments;
    \item attach previously unseen sensors;
    \item expose previously unseen actuators;
    \item provide long temporal dependencies;
    \item create conflicting and later reversed regularities;
    \item expose text, audio, visual contact, action, and consequences together;
    \item require learned lower relations to become material for higher relations;
    \item require higher relations to affect prediction and action;
    \item require correction to propagate while preserving independently supported structure.
\end{enumerate}

The expected developmental chain is:

\[
\boxed{
\text{contact}
\rightarrow
\text{relations}
\rightarrow
\text{reusable abstractions}
\rightarrow
\text{relations over relations}
\rightarrow
\text{prediction/action}
\rightarrow
\text{new reality contact}
\rightarrow
\text{correction and further growth}
}
\]

This chain must continue under newly introduced domains without adding domain-specific cognitive authority.

\section{Final Blueprint}

The complete architecture reduces to:

\[
\boxed{
\begin{aligned}
M_{t+1} &= C(M_t\oplus R_t),\\[4pt]
R_t &= \Sx(\W_t),\\[4pt]
\W_{t+1} &\supseteq \Ax(M_{t+1}),\\[4pt]
\text{Memory} &= M,\\[4pt]
\text{Relations} &\text{ are first-class and recursively reusable},\\[4pt]
\text{Development} &=
\text{continuous reality-constrained relational growth}.
\end{aligned}
}
\]

In words:

\begin{quote}
One persistent mathematical individual continuously receives reality contact, preserves it in a single relational whole, forms relations over anything already present in that whole, recursively reuses those relations as new material, corrects them when reality disagrees, and allows the resulting structure to alter the same individual's future perception, prediction, expression, and action.
\end{quote}

\end{document}
